% ============================================================
%  Helge Holst, "The Principle of Relativity"
%  Fysisk Tidsskrift, 15. Aargang (1916-17), pp. 165-168.
%
%  ENGLISH TRANSLATION
%  Status: COMPLETE. Translated from transcription.tex (the Danish diplomatic
%  transcription, COMPLETE), never from the scan. Structure mirrors
%  transcription.tex 1:1: same paragraphing, same page markers, the one
%  footnote at the same anchor, all 8 emphasis runs carried over.
%
%  Holst's SECOND contribution to the Kroman relativity debate -- the sixth of
%  its seven items. Not to be confused with his »Tidsproblemet« at pp. 1-13 of
%  the same volume, which is ../holst-tidsproblemet. It answers H. M. Hansen's
%  article at pp. 18-29 of this volume (../hansen-svar), which remains under
%  the repo's embargo until 1 January 2027; Holst's own text is free.
%  See ../relativitetsprincipet/note.md for the debate as a whole.
%
%  RIGHTS: clear. Helge Holst died in 1944; Danish copyright (life+70)
%  expired at the end of 2014.
%
% ------------------------------------------------------------
%  TRANSLATION CONVENTIONS -- this article
% ------------------------------------------------------------
%  * Standing method: ../../../TRANSLATION-PLAYBOOK.md; the cluster's shared
%    decisions are set out at length in
%    ../relativitetsprincipet/translation.tex.
%  * EMPHASIS. All 8 \emph{} runs carried over. As in »Tidsproblemet«, and
%    unlike the Kroman articles, PROPER NAMES ARE NEVER LETTERSPACED here --
%    Lorentz, Einstein, Michelson all plain. The English follows; the absence
%    of \emph{} on the names is not an oversight.
%  * TITLE SPELLING. The Danish display title and running heads read
%    »Relativitetsprincipet« with ONE p, while the body writes
%    »Relativitetsprincippet« with TWO. The distinction cannot be carried into
%    English and is recorded here instead; transcription.tex keeps both.
%  * "Dr. H." is Holst's abbreviation for H. M. Hansen throughout, and "M.'s
%    Forsøg" on p. 167 is Michelson's experiment. Both abbreviations are kept
%    as he sets them rather than expanded.
%  * PRINTER'S DEFECTS. p. 166's missing comma pair round the parenthetical
%    "in following the rules of synchronisation" is reproduced. p. 167's »at«
%    for »af« breaks the syntax of the sentence; the English reproduces the
%    broken construction rather than repairing it, and a comment marks it.
%
%  BUILD: pdflatex (libertinus + libertinust1math), like the Danish.
% ============================================================

\documentclass[12pt,a4paper]{article}
\usepackage[utf8]{inputenc}
\usepackage[T1]{fontenc}
\usepackage{libertinus}
\usepackage{libertinust1math}
\usepackage{textalpha}
\usepackage{amsmath}

% FOOTNOTES: exactly one, on p. 166, marked *). A fixed mark is correct here.
\renewcommand{\thefootnote}{*)}
\usepackage[english]{babel}
\usepackage{enumitem}
\usepackage[protrusion=true,expansion=true]{microtype}
\usepackage{setspace}
\usepackage[margin=1.2in]{geometry}
\usepackage{fancyhdr}
\setlength{\headheight}{15pt}
\usepackage{hyperref}

\hypersetup{
  pdftitle={The Principle of Relativity (Holst 1917)},
  pdfauthor={Helge Holst},
  hidelinks
}

\pagestyle{fancy}
\fancyhf{}
\fancyhead[LE,RO]{\thepage}
\fancyhead[RE]{\textit{The Principle of Relativity}}
\fancyhead[LO]{\textit{Helge Holst}}

\setstretch{1.3}

\title{\textbf{``The Principle of Relativity''}}
\author{H.\ Holst}
\date{\textit{Fysisk Tidsskrift} 15 (1916--17), pp.~165--168\\[2pt]
\small Translated from the Danish}

\begin{document}
\maketitle

% --- p. 165 ---
% The article BEGINS ON p. 165, in the lower half of the page, under the
% running head of the preceding article. Title block supplied by \maketitle.
On the occasion of Dr.~H.~M. Hansen's article on the principle of relativity
in the 1st number of this volume, I should like to ask for space for a few
more remarks, in the hope that these may contribute to clarifying further the
relation between the different standpoints towards Einstein's doctrine.

In so far as one merely applies the Einsteinian \emph{procedure} purely
phenomenologically, as a means of eliminating apparent contradictions (for
example those laid bare by the Michelson experiment), no epistemological
objections can be raised against it. One may of course adopt the rule that
each single one of several observers in different systems, which move in
relation to one another, acts as though light in his system always has the
constant velocity $c$ in all directions (and that he therefore synchronises
his clocks by light signals in the way indicated by Einstein), and further
that one passes over from the one of the systems to the other by the
Einsteinian transformation formulae. One has further the right to set up the
hypothesis that the results one finds under the application of these rules
will in all domains agree with the experimental re\-%
% --- p. 166 ---
sults. If it now turns out that one obtains such agreement not merely for the
phenomena for which the procedure was developed, but also for quite other
phenomena, indeed that one can predict new experimental results which are
afterwards really found, then the hypothesis has established itself as a
fruitful guiding hypothesis, and it is understandable that those who have
worked with it become inclined to rely on its correctness and to regard it as
an expression of fundamental peculiarities of the world. The circumstance that
the different observers in following the rules of synchronisation come into
% PRINTER'S DEFECT, as printed: the parenthetical "in following the rules of
% synchronisation" is unpunctuated in the Danish -- no comma before it and
% none after. Reproduced here rather than repaired.
disagreement about determinations of simultaneity, and that the hypothesis
entails that no \emph{criterion} can be found as to who is right, does not in
and for itself make it epistemologically unacceptable that the hypothesis of
the procedure's general applicability should be correct. Lorentz has of course
by his hypothesis\footnote{Both here and in the earlier article there is meant
by Lorentz's hypothesis or interpretation not merely the original contraction
hypothesis, but this together with its later additions.} shown that it is (at
all events apparently) \emph{possible} to give a physical explanation of the
matter completely reconcilable with the natural conception of space and time.

But as Einstein's doctrine has been put forward by its originator, and as it
is generally presented, it is not merely a description of a procedure and the
setting up of a hypothesis (or a dogma) about that procedure's general
applicability. Besides a phenomenological side it has also an epistemological
one, in that it is expressly maintained or tacitly presupposed that the rules
for determinations of simultaneity are more than purely conventional, that the
velocity of light is not merely \emph{set as} constant but \emph{is} constant
and the same in relation to sender and receiver, even if these move in
relation to one another, that no \emph{meaning} can be attached to speaking of
absolute simultaneity, and that a physical explanation of the contradictions
mentioned becomes superfluous when one revolutionises one's conception of
space and time.

It is against this side of the matter that protests have arisen. I for my part
have sought to show that our thought really does get ``a place where it can
stand'', and from which it can view Einstein's doctrine from outside, provided
only that one maintains its right to increase velocities beyond all limits. If
we gave up this right for thought, because there is (possibly) a definite
physical limit to attainable velocities, then we should for example just as
well have to give up thought's right to
% --- p. 167 ---
prolong straight lines beyond all limits, if some well-founded physical theory
had led to the conclusion that our world-system was enclosed in a completely
impenetrable shell. Fortunately for physics, its indispensable helper
mathematics makes itself no physical scruples of this kind; it is one of its
inalienable rights to let thought go beyond the physically attainable. But
when thought does not let itself be stopped by physical barriers, we can, as
is developed in my article, attach a definite meaning to (and give a
definition of) two events being really simultaneous, even if the
world-mechanism should be so arranged that we cannot find any criterion of it.
% »Kriterium« is NOT letterspaced here in the Danish, though it is on p. 166.
% The English follows: no \emph{} on this occurrence.

Dr.~H. acknowledges (p.~28, l.~13 from the foot) that this conception of the
concept of simultaneity (which coincides not merely with the one indicated by
Lorentz, but with all men's tacit presuppositions --- until the year 1905) can
be maintained; but that the underlining of the word \emph{can}, as well as
% PRINTER'S DEFECT, as printed: the Danish has »men at Understregningen ...
% samt af andre Ytringer ... synes det at fremgaa«, where the sense wants
% »men af«. The »at« leaves the sentence without a subject for »synes det at
% fremgaa«. The English reproduces the broken construction rather than
% mending it.
from other utterances in the article, it seems to appear that he thinks it is
not necessary for thought to maintain it. In this I disagree with him. The
conception named can and must be maintained as the one which follows from our
natural concepts of space and time and from the fundamental laws of our
thinking. It is, I think, an illusion that we can get away from this, and it
is therefore also an illusion that Einstein's doctrine gives any contribution
whatever to the \emph{explanation} of the Michelson difficulties; it merely
gives simple rules for how, by the introduction of certain factors, one can
eliminate the disagreements which arise, for example, between physicists who
on the earth and on Mars carry out M.'s experiment.

It is another matter that when one applies Einstein's procedure purely
phenomenologically, the whole question of the conception of the concept of
simultaneity goes out of the discussion, and it may be said to be without
interest for the working physicist, in so far as he can for the present push
the matter aside.

Undoubtedly Dr.~H.'s full and instructive expositions of Einstein's doctrine
have been of much benefit, and there is reason to be grateful to him for the
liberality with which he has in the journal given space for contributions
which go against his own points of view. When in his last article he stresses
the phenomenological so strongly, he at the same time opens a good way out
towards the levelling of the oppositions, and a prospect that all physicists
may with pleasure follow a fruitful application of Ein\-%
% --- p. 168 ---
stein's method. This is made possible in the higher degree, the more
completely it succeeds in cleansing modes of expression and lines of
consideration of the epistemological elements which set thought reacting.

One must however hope that the physicists who work with Einstein's procedures
may not be so absorbed in the work --- however fine and great the results it
might yield --- that they wholly lose their interest in physical explanations
of the difficulties, or settle so firmly into belief in the Einsteinian
hypothesis that they reject in advance attempts at explanation which, in their
further development (in contrast to Lorentz's), might lead to the Einsteinian
transformation equations being able to give only approximately correct
results.
% The article ENDS HERE, about a third of the way down p. 168. No signature
% and no dateline; a new item begins on the rest of the page.

\end{document}
